\documentclass[12pt]{article}
\usepackage[utf8]{inputenc}
\usepackage[T1]{fontenc}
\usepackage[margin=1in]{geometry}
\usepackage{setspace}
\usepackage{xcolor}
\usepackage[colorlinks=true, urlcolor=blue!60!black, linkcolor=blue!60!black]{hyperref}
\setlength{\parindent}{2em}
\setlength{\parskip}{1.25em}
\renewcommand{\baselinestretch}{1.0}
\title{\textit{NLRB v. Noel Canning}, 573 U.S. 513 (2014)}
\author{Case Brief}
\date{}
\begin{document}
\maketitle

\section{Case Metadata}
\begin{itemize}
  \item \textbf{Case Name:} \emph{National Labor Relations Board v. Noel Canning}, \href{https://supreme.justia.com/cases/federal/us/573/513/}{573 U.S. 513 (2014)}
  \item \textbf{Docket No.:} 12-1281 \quad \textbf{Argued:} January 13, 2014 \quad \textbf{Decided:} June 26, 2014
  \item \textbf{Court:} Supreme Court of the United States (on writ of certiorari to the U.S. Court of Appeals for the District of Columbia Circuit, 705 F.3d 490 (D.C. Cir. 2013))
  \item \textbf{Briefing Book Placement:} Chapter 2: Separation of Powers and the Executive
  \item \textbf{Vote:} 9--0 judgment affirming the D.C. Circuit (5--4 split on textual scope of the Recess Appointments Clause)
  \item \textbf{Opinion of the Court:} Justice Breyer (joined by Justices Kennedy, Ginsburg, Sotomayor, and Kagan)
  \item \textbf{Concurring in the Judgment:} Justice Scalia (joined by Chief Justice Roberts and Justices Thomas and Alito)
  \item \textbf{Dissenting Opinions:} None
  \item \textbf{Key Constitutional Provisions:} Appointments Clause, U.S. Const. art. II, \S~2, cl. 2; Recess Appointments Clause, art. II, \S~2, cl. 3; Rules of Proceedings Clause, art. I, \S~5, cl. 2; Adjournments Clause, art. I, \S~5, cl. 4; Presidential Adjournment Clause, art. II, \S~3
  \item \textbf{Primary Source Links:} \href{https://www.law.cornell.edu/supremecourt/text/12-1281}{Cornell LII Supreme Court Collection}; \href{https://supreme.justia.com/cases/federal/us/573/513/}{Justia U.S. Supreme Court Center}
\end{itemize}

\section{Facts and Procedural History}
Noel Canning, a Pepsi-Cola bottler and distributor in Yakima, Washington, was found by the National Labor Relations Board (NLRB) to have unlawfully refused to reduce to writing and execute a collective-bargaining agreement with a labor union (\emph{Noel Canning}, 358 N.L.R.B. No. 4 (2012)). Noel Canning petitioned the D.C. Circuit to set aside the Board's order, contending that the five-member Board lacked the three lawfully appointed members required for a statutory quorum under 29 U.S.C. \S~153(b) and \emph{New Process Steel, L.P. v. NLRB}, 560 U.S. 674 (2010), because three of its members---Sharon Block, Richard Griffin, and Terence Flynn---had been invalidly appointed (\href{https://www.law.cornell.edu/supremecourt/text/12-1281}{Cornell LII}).

President Obama had nominated all three individuals in 2011 (Flynn's nomination had been pending roughly a year; Block's and Griffin's for several weeks). On December 17, 2011, the Senate adopted by unanimous consent a resolution calling for brief ``pro forma session[s],'' with ``no business . . . transacted,'' every Tuesday and Friday from December 20, 2011, through January 20, 2012, before resuming ordinary business on January 23, 2012. On January 4, 2012---during the three-day break between the January 3 and January 6 pro forma sessions---President Obama invoked the Recess Appointments Clause to appoint Block, Griffin, and Flynn to the NLRB (\href{https://www.law.cornell.edu/supremecourt/text/12-1281}{Cornell LII}).

The D.C. Circuit vacated the NLRB order, holding on narrow textualist grounds that ``the Recess of the Senate'' covers only \emph{inter-session} recesses (breaks between formal annual sessions) and that ``Vacancies that may happen during the Recess'' covers only vacancies that \emph{first arise} during that recess (705 F.3d 490 (D.C. Cir. 2013)). The Supreme Court granted certiorari (570 U.S. 916 (2013)), adding the third question of whether the recess-appointment power may be exercised while the Senate convenes every three days in pro forma sessions, and unanimously affirmed the judgment on that third ground.

\section{Issue(s) Presented}
\begin{enumerate}
  \item Does ``the Recess of the Senate'' in Article II, \S~2, cl. 3 encompass both inter-session and intra-session recesses, and how long must a recess be to trigger the Clause?
  \item Does ``Vacancies that may happen during the Recess of the Senate'' include vacancies that arose prior to a recess and continue to exist during it, or only vacancies that initially occur during the recess?
  \item For purposes of calculating the length of a recess, must the Senate's pro forma sessions be treated as sessions, or ignored as functional recesses?
\end{enumerate}

\section{Holding}
\begin{enumerate}
  \item \textbf{Both inter-session and intra-session recesses (if of sufficient length).} The Clause applies to both kinds of recesses, but a break of three days or fewer is categorically too short, and a recess of fewer than 10 days is \emph{presumptively} too short absent very unusual circumstances (\href{https://www.law.cornell.edu/supremecourt/text/12-1281}{Cornell LII}).
  \item \textbf{Both pre-existing and recess-arising vacancies.} The Clause empowers the President to fill any vacancy that exists during a qualifying recess, even if arose before the recess began.
  \item \textbf{Pro forma sessions count as sessions.} The Senate is in session when it says it is, provided that under its own rules it retains the capacity to transact Senate business. Because the Senate was in session on January 3 and January 6, 2012, the intervening three-day break was too short to trigger the Clause, rendering President Obama's January 4, 2012 NLRB appointments unconstitutional.
\end{enumerate}

\section{Reasoning}
Writing for a five-Justice majority, Justice Breyer framed the Clause as a ``subsidiary'' mechanism designed to ensure executive continuity while the Senate is unavailable, and anchored his interpretive method in \textbf{historical practice}---invoking James Madison's concept of constitutional ``liquidation'' (\emph{McCulloch v. Maryland}; \emph{The Pocket Veto Case}; Justice Frankfurter's \emph{Youngstown} concurrence) to give ``great weight'' to longstanding interbranch working arrangements when constitutional text is ambiguous (\href{https://www.law.cornell.edu/supremecourt/text/12-1281}{Cornell LII}).

\textbf{1. Inter-Session and Intra-Session Recesses (Part III).} Breyer found the phrase ``the Recess'' linguistically ambiguous: Founding-era dictionaries and usages employed ``the recess'' generically for intra-session breaks as well as formal inter-session intervals. Functionally, the Senate is equally unavailable to give advice and consent during a lengthy intra-session break. Historically, Presidents since Andrew Johnson (1867--1868)---and routinely since 1921 and World War II---have made thousands of intra-session recess appointments, backed by nearly unanimous Executive Branch legal opinions and unopposed by any formal Senate action (indeed, the 1905 Senate Judiciary Committee Report defined ``recess'' functionally). Yet the Adjournments Clause (Art. I, \S~5, cl. 4) treats breaks of three days or fewer as constitutionally \emph{de minimis}, and two centuries of practice revealed zero intra-session recess appointments during breaks shorter than 10 days. Thus, recesses under 10 days are presumptively too short, leaving open only extraordinary emergencies such as a national catastrophe (explicitly excluding Senate political opposition) (\href{https://www.law.cornell.edu/supremecourt/text/12-1281}{Cornell LII}).

\textbf{2. Pre-Existing Vacancies (Part IV).} Although ``happen'' most naturally means ``initially occur'' to a modern ear, Breyer noted Founding-era ambiguity---citing Thomas Jefferson's 1802 observation that it could mean ``happen to be'' and Attorney General William Wirt's 1823 opinion reading it as ``happen to exist.'' The narrower reading would paralyze essential government offices whenever a vacancy arose shortly before adjournment. Two centuries of practice since President Madison, near-unanimous Attorney General opinions, and Congress's 1940 Pay Act amendments (5 U.S.C. \S~5503, permitting salary payments for certain pre-recess vacancies) confirmed the broader view (\href{https://www.law.cornell.edu/supremecourt/text/12-1281}{Cornell LII}).

\textbf{3. Pro Forma Sessions (Part V).} Turning to the January 2012 appointments, Breyer refused the Solicitor General's invitation to pierce the Senate's formal record via C-SPAN footage. Under the Rules of Proceedings Clause (Art. I, \S~5, cl. 2), \emph{United States v. Ballin} (1892), and \emph{Marshall Field \& Co. v. Clark} (1892), courts must respect the Senate's own declaration that it is in session so long as it retains the \emph{capacity} under its rules to transact business. Because Senate rules presume a quorum unless questioned and allow legislation or confirmations by unanimous consent---as the Senate actually did when it enacted a payroll-tax extension during its December 23, 2011 pro forma session (Pub. L. 112-78)---the pro forma sessions broke the holiday adjournment into three-day intervals, defeating the January 4 appointments (\href{https://www.law.cornell.edu/supremecourt/text/12-1281}{Cornell LII}).

\section{Concurring and Dissenting Opinions}
\begin{itemize}
  \item \textbf{Justice Scalia (joined by Chief Justice Roberts and Justices Thomas and Alito), concurring in the judgment:} Agreed unanimously that the 2012 NLRB appointments were void, but fiercely attacked the majority's textual and methodological premises. Scalia argued that structural separation-of-powers rules protect individual liberty, not branch convenience, and cannot be altered by interbranch acquiescence (\emph{INS v. Chadha}; \emph{Free Enterprise Fund}). On original public meaning, ``the Recess'' is a formal term paired with ``the next Session,'' meaning solely the inter-session break, and ``Vacancies that may happen'' unambiguously means vacancies that \emph{arise} during that recess---confirmed by Attorneys General Randolph and Lee, the First Congress's statutes, and the 1863 Pay Act. Scalia ridiculed the majority's 10-day presumption as unmoored judicial legislation and warned that requiring formal Senate-wide opposition to block executive ``adverse possession'' structurally favors presidential aggrandizement (\href{https://www.law.cornell.edu/supremecourt/text/12-1281}{Cornell LII}).
  \item \textbf{Dissents:} None.
\end{itemize}

\section{Key Quotations}
\begin{quote}
  ``We have not previously interpreted the Clause, and, when doing so for the first time in more than 200 years, we must hesitate to upset the compromises and working arrangements that the elected branches of Government themselves have reached.''
  --- Justice Breyer, \href{https://www.law.cornell.edu/supremecourt/text/12-1281}{573 U.S. at 526 (Cornell LII)}.
\end{quote}
\begin{quote}
  ``In sum, we conclude that the phrase `the recess' applies to both intra-session and inter-session recesses. If a Senate recess is so short that it does not require the consent of the House, it is too short to trigger the Recess Appointments Clause. See Art. I, \S~5, cl. 4. And a recess lasting less than 10 days is presumptively too short as well.''
  --- Justice Breyer, \href{https://www.law.cornell.edu/supremecourt/text/12-1281}{573 U.S. at 538 (Cornell LII)}.
\end{quote}
\begin{quote}
  ``We hold that, for purposes of the Recess Appointments Clause, the Senate is in session when it says it is, provided that, under its own rules, it retains the capacity to transact Senate business.''
  --- Justice Breyer, \href{https://www.law.cornell.edu/supremecourt/text/12-1281}{573 U.S. at 550 (Cornell LII)}.
\end{quote}
\begin{quote}
  ``The majority justifies those atextual results on an adverse-possession theory of executive authority: Presidents have long claimed the powers in question, and the Senate has not disputed those claims with sufficient vigor, so the Court should not `upset the compromises and working arrangements that the elected branches of Government themselves have reached.' ''
  --- Justice Scalia, concurring in the judgment, \href{https://www.law.cornell.edu/supremecourt/text/12-1281}{573 U.S. at 570--71 (Cornell LII)}.
\end{quote}

\section{Doctrinal Significance and Modern Legacy (Through 2026)}
\emph{Noel Canning} has exerted a dual legacy across separation-of-powers doctrine through 2026. Practically, it handed the Senate a near-airtight procedural shield: by holding twice-weekly pro forma sessions, a Senate controlled by the opposition party---or even a single chamber refusing consent to a >3-day adjournment---can block recess appointments indefinitely. Consequently, Presidents after 2014 largely shifted to temporary acting designations under the Federal Vacancies Reform Act (FVRA), whose statutory limits the Court strictly policed in \href{https://supreme.justia.com/cases/federal/us/580/15-1251/}{\emph{NLRB v. SW General, Inc.}, 580 U.S. 288 (2017)} (distinguishing \emph{Noel Canning}'s long historical gloss). Following the 2024 election, however, renewed debate erupted over bypassing pro forma sessions---including proposals for a friendly Senate to take a 10-day recess or for the President to invoke the never-used Article II, \S~3 Disagreement/Adjournment Clause (cited in \emph{Noel Canning}, 573 U.S. at 550--51) to force a recess (\href{https://reason.com/volokh/2024/11/10/president-elect-trumps-tweet-about-recess-appointments/}{Volokh Conspiracy, Nov. 10, 2024}; \href{https://americarenewing.com/issues/brief-on-the-article-ii-recess-appointments-clause/}{Center for Renewing America, Nov. 17, 2024}).

Methodologically, \emph{Noel Canning} became the modern touchstone for \textbf{historical practice and constitutional ``liquidation''} in separation-of-powers litigation. The Court has repeatedly invoked \emph{Noel Canning}'s liquidation framework---both to validate longstanding interbranch norms in \href{https://supreme.justia.com/cases/federal/us/576/1/}{\emph{Zivotofsky v. Kerry}, 576 U.S. 1 (2015)} (exclusive presidential recognition power), \href{https://supreme.justia.com/cases/federal/us/591/19-715/}{\emph{Trump v. Mazars USA, LLP}, 591 U.S. 848 (2020)} (congressional subpoenas for presidential papers), and \href{https://supreme.justia.com/cases/federal/us/601/22-448/}{\emph{CFPB v. Community Financial Services Ass'n}, 601 U.S. 416 (2024)} (historical appropriations practice), and conversely to strike down novel executive structures lacking historical pedigree in \emph{Seila Law LLC v. CFPB} (2020) and \href{https://supreme.justia.com/cases/federal/us/609/25-332/}{\emph{Trump v. Slaughter}, 609 U.S. \_\_\_ (2026)} (relying on the Decision of 1789 and early presidential practice).

\section{Pedagogical Critiques and Exam Hypotheticals}
\textbf{Pedagogical Critiques.} First, \emph{Noel Canning} showcases a striking methodological paradox: Justice Breyer adopts a sweeping functionalist/historical-gloss approach on Questions 1 and 2 (rejecting formalistic distinctions between inter- and intra-session recesses), yet pivots to strict formalism on Question 3 (accepting nominal 30-second pro forma sessions at face value). Second, Scalia's institutional critique of ``adverse possession'' remains potent in constitutional theory: because the unitary President acts decisively while the multimember Senate faces severe collective-action hurdles, treating legislative silence as ``acquiescence'' creates a one-way ratchet favoring executive expansion. Third, the 10-day presumption illustrates the tension between bright-line administrability and textual fidelity---neither ``3 days'' nor ``10 days'' appears in Article II, \S~2, cl. 3.

\textbf{Exam Hypothetical and ``Opposite Side'' Contrast.}
\begin{itemize}
  \item \textbf{Scenario A (Valid Recess Appointment Under \emph{Noel Canning}):} A Cabinet Secretary resigns on June 1 while the Senate is in session. On August 1, both Houses of Congress adopt a concurrent resolution adjourning for an 18-day intra-session summer recess with no pro forma sessions scheduled. On August 5, the President recess-appoints a nominee whom the Senate had previously declined to advance out of committee. \emph{Analysis:} Valid under \emph{Noel Canning}. Even though the vacancy pre-existed the recess and the break is intra-session, the recess exceeds 10 days, the Senate has formally adjourned without pro forma sessions, and the appointee may serve until the end of the Senate's next formal session (subject to Pay Act limits under 5 U.S.C. \S~5503).
  \item \textbf{Scenario B (``Opposite Side'' Contrast---Pro Forma Sessions \& Short Break):} Now suppose that instead of adjourning for 18 days, the Senate convenes every Tuesday and Friday in 45-second pro forma sessions presided over by a single Senator, declaring no business will be conducted, while the House has adjourned for August. Frustrated by partisan gridlock, the President makes a recess appointment on a Thursday (during a 3-day gap) and alternatively argues that an 8-day holiday break without pro forma sessions qualifies because Senate obstruction is an ``unusual circumstance.'' \emph{Analysis:} Invalid on both counts under \emph{Noel Canning}. Under Part V, the Senate is in session during its pro forma sessions because it retains capacity to act by unanimous consent, making the 3-day gap categorically too short. And under Part III, an 8-day break is presumptively too short, and the Court explicitly held that ``political opposition in the Senate'' never qualifies as an unusual circumstance overcoming the 10-day presumption.
\end{itemize}

\end{document}
